\chapter{G\.{I}R\.{I}\c{S}}

\label{giris} 
%%%%%%%%%%%%%%%%%%%%%%%%%%%%%%%%%%%%%%%%%%%%%%%%%%%%%%%%%%%%%%%%%

Giri\c{s} b\"{o}l\"{u}m\"{u}, tez \c{c}al\i{}\c{s}mas\i{}n\i{}n temel \c{c}er%
\c{c}evesini okuyucuya sunar. Bu b\"{o}l\"{u}mde ara\c{s}t\i{}rma konusu,
problem durumu, tezin amac\i{} ve \"{o}nemi, s\i{}n\i{}rl\i{}l\i{}klar,
varsay\i{}mlar ve temel tan\i{}mlar akademik bir b\"{u}t\"{u}nl\"{u}k i\c{c}%
inde a\c{c}\i{}klanmal\i{}d\i{}r.

K\i{}lavuz gere\u{g}i G\.{I}R\.{I}\c{S} ana ba\c{s}l\i{}\u{g}\i{} alt\i{}nda
ayr\i{}ca numaral\i{} alt ba\c{s}l\i{}k olu\c{s}turulmamal\i{}d\i{}r. Bu
nedenle bu \c{s}ablonda Giri\c{s} b\"{o}l\"{u}m\"{u}, alt ba\c{s}l\i{}k
kullanmadan tek bir ak\i{}\c{s} halinde verilmi\c{s}tir. Gerekli bilgiler,
paragraflar aras\i{}nda mant\i{}kl\i{} ge\c{c}i\c{s}ler kurularak sunulmal\i{%
}; y\"{o}ntem, literat\"{u}r ve bulgulara ili\c{s}kin ayr\i{}nt\i{}lar
ilgili ana b\"{o}l\"{u}mlere b\i{}rak\i{}lmal\i{}d\i{}r.

Bu b\"{o}l\"{u}mde okuyucuya \c{c}al\i{}\c{s}man\i{}n hangi bilimsel bo\c{s}%
lu\u{g}a odakland\i{}\u{g}\i{}, hangi sorulara yan\i{}t arad\i{}\u{g}\i{} ve
tezin genel organizasyonunun nas\i{}l ilerledi\u{g}i k\i{}sa ve a\c{c}\i{}k
bi\c{c}imde aktar\i{}labilir.
